\documentclass[11pt]{article}
\usepackage{fontspec}
\usepackage[french]{babel}
\usepackage{amsmath}
\usepackage{amssymb}
\usepackage{amsthm}
\usepackage{mathrsfs}
\usepackage{tikz-cd}
\usepackage{enumitem}
\usetikzlibrary{babel}

\newcommand{\printedpage}[1]{\marginpar{\footnotesize p.~#1}}
\newcommand{\msp}[1]{\marginpar{\footnotesize [#1]}}
\newcommand{\D}{\mathbb{D}}
\newcommand{\Dia}{\mathsf{Dia}}
\newcommand{\Ob}{\mathsf{Ob}}
\newcommand{\Fl}{\mathsf{Fl}}
\newcommand{\Hom}{\underline{\mathsf{Hom}}}
\newcommand{\hot}{\mathsf{hot}}
\newcommand{\Cat}{\mathcal{C}at}
\newcommand{\isoto}{\overset{\sim}{\longrightarrow}}

\begin{document}

\printedpage{39}
% head: A. Grothendieck : Dérivateurs, Ch. I, Édition des Universités de Montpellier II et Paris VII
% foot: Version 8 juillet 2001
\msp{page 67}

\noindent Donc la flèche envisagée $(**)$ s'interprète comme
\[
H_*^{\D}(y\backslash X, M_X) \longrightarrow M
\]
et c'est l'homomorphisme d'adjonction $p_!(p^*(M)) \to M$.
\[
\begin{tikzcd}
X \arrow[rr, "f"] \arrow[dr, "p"{swap}] & & Y \arrow[dl, "q"] \\
& e &
\end{tikzcd}
\]

\noindent Donc c'est un isomorphisme pour tout $M$ si et seulement si $y\backslash X$ est $\D$-contractile.

\noindent On va énoncer ce qu'on a prouvé.

\noindent\textbf{Théorème 1.} \textit{Soit}
\[
\begin{tikzcd}
X \arrow[rr, "f"] \arrow[dr, "p"{swap}] & & Y \arrow[dl, "q"] \\
& e &
\end{tikzcd}
\]
\textit{dans} $\Dia$. \textit{On suppose que} $\D$ \textit{est un prédérivateur satisfaisant} \textbf{Der 2}, \textit{que} $p_*, q_*$ \textit{et} $f_!$ \textit{existent, que} $f_!$ \textit{satisfait à} \textbf{Der 4} (d) \textit{pour le calcul des fibres. Conditions équivalentes :}

a) $f$ \textit{est une} $\D$\textit{-équivalence costricte, i.e.} (prop. 2) \textit{pour tout} $y \in \Ob\, Y$, $y\backslash X$ \textit{est} $\D$\textit{-contractile ;}

b) \textit{pour tout} $F \in \Ob\, \D(Y)$, \textit{l'homomorphisme suivant est} \texttt{[un]} \textit{isomorphisme}
\[
H^*_{\D}(Y, F) \isoto H^*_{\D}(X, f^*(F)) \ ;
\]

c) \textit{pour tout} $M \in \Ob(\mathcal{A} = \D(e))$, \textit{l'homomorphisme d'adjonction est} \texttt{[un]} \textit{isomorphisme}
\[
f_!(M_X) = f_!(f^*(M_Y)) \isoto M_Y \ .
\]

Je vais renuméroter les axiomes sur le prédérivateurs. % [sic] le prédérivateurs -> les prédérivateurs, article singulier devant un pluriel

\noindent\textbf{Der 1}, \textbf{Der 2} comme avant (p. 45).

\noindent\textbf{Der 3}\footnote[28]{Axiome des catégories d'indices asphériques, \textit{cf.} variantes \textbf{Der 3} a) b) c) page 97, et surtout pages 111-112.} : Soit $X$ dans $\Dia$, $p : X \to e$. Supposons $X$ contractile.\footnote[29]{N.B. Il suffit de supposer que $X$ a un objet final. Mais ce n'est pas un axiome autodual.} Alors
\[
p^* : M \longmapsto M_X \quad , \qquad p^* : \mathcal{A} = \D(e) \longrightarrow \D(X)
\]
est pleinement fidèle.

\printedpage{40}
% head: A. Grothendieck : Dérivateurs, Ch. I, Édition des Universités de Montpellier II et Paris VII
% foot: Version 8 juillet 2001
\msp{page 68}

\noindent\textbf{Der 4} (g) : existence des $f_*$.

\noindent\textbf{Der 4} (d) : existence des $f_!$.

\noindent\textbf{Der 5} (g) Calcul des fibres de $f_*(F)$.

\noindent\textbf{Der 5} (d) Calcul des fibres de $f_!(F)$.

\noindent L'ancien axiome \textbf{Der 5} (g) résulte de \textbf{Der 2}, \textbf{Der 3}, \textbf{Der 4} (g,d), \textbf{Der 5} (d), grâce au théorème précédent.

Un aspect intéressant du théorème précédent, si on l'envisage comme énonçant l'équivalence de a) et b), c'est que cet énoncé ne fait pas appel à des notions homologiques, mais que sa démonstration y fait appel, sous une forme relativement forte (calcul des fibres de $f_!(F)$). Pour les théorèmes de changement de base pour la $\D$-cohomologie relative, \textit{i.e.} pour les $f_*^{\D}$, on a besoin aussi constamment du calcul des fibres des $f_*(F)$. Donc il semble qu'on ne peut se passer d'aucun des axiomes précédents, y compris \textbf{Der 4} (d) et \textbf{Der 5} (d), même si on ne veut s'intéresser qu'au formalisme des $f^*, f_*$ (aspect cohomologique). Par la suite, on va donc supposer ces axiomes satisfaits.

\noindent\textbf{Théorème 2.} \textit{Soit} $f : X \to Y$ \textit{dans} $\Dia$. \textit{Conditions équivalentes :}

a) \textit{Pour tout} $y_0 \in \Ob\, Y$, \textit{le foncteur canonique}
\[
X_{y_0} \overset{\varphi}{\longrightarrow} X/y_0 \quad (\simeq X \times_Y Y/y_0) \ ,
\]
\textit{où} $X_{y_0}$ \textit{désigne la fibre en} $y_0$, \textit{est une} $\D$\textit{-équivalence co-stricte.}

b) \textit{Pour tout} $x_0 \in \Ob\, X$, \textit{le foncteur induit sur les catégories colocalisées en} $x_0$ \textit{et en} $y_0 = f(x_0)$
\[
x_0\backslash X \longrightarrow y_0\backslash Y
\]
\textit{est à fibres} $\D$\textit{-contractiles.}

\msp{page 69}

c) \textit{La formation des} $f_*F$ \textit{(}$F$ \textit{dans} $\D(X)$\textit{) commute à tout changement de base} $Y' \to Y$, \textit{et ça reste vrai après tout changement de base -- i.e. si on a} $Y'' \to Y' \to Y$, \textit{d'où}
\[
\begin{tikzcd}
X'' \arrow[r] \arrow[d, "f''"{swap}] & X' \arrow[r] \arrow[d, "f'"] & X \arrow[d, "f"] \\
Y'' \arrow[r, "p'"{swap}] & Y' \arrow[r, "p"{swap}] & Y
\end{tikzcd}
\]
\textit{le “théorème de changement de base” est vrai pour} $f'_*$ \textit{et} $p'^{*}$.

c$'$) \textit{La formation des} $f_*(F)$ \textit{(}$F \in \Ob\, \D(X)$\textit{) commute au passage aux fibres et ceci reste vrai après tout changement de base} $p : Y' \to Y$.

c$''$) \textit{Comme} c$'$) \textit{avec} $Y' \to Y$ \textit{un morphisme de localisation} $Y/y \to Y$.

d) \textit{Si} $Y'' \overset{p}{\to} Y' \to Y$, \textit{avec} $p$ \textit{coasphérique, alors} $q : X'' = X \times_Y Y'' \to X' = X \times_Y Y'$ \textit{est coasphérique.}

\printedpage{41}
% head: A. Grothendieck : Dérivateurs, Ch. I, Édition des Universités de Montpellier II et Paris VII
% foot: Version 8 juillet 2001

d$'$) \textit{Comme} d) \textit{avec} $Y'' = e$ (donc $Y'$ est catégorie avec objet final $e$).

d$''$) \textit{Comme} d) \textit{avec} $Y'' = \{1\}$, $Y' = \boldsymbol{\Delta}_1$.

.

\noindent\textsc{Démonstration}. a) $\Leftrightarrow$ b) cf. proposition page 53. Montrons que ces conditions impliquent c). Comme elles sont stables par changement de base (\textit{cf.} démonstration immédiate page 55) donc par \textbf{Der 2} on est ramené à prouver que $H^*_{\D}(X/y, F) \overset{\sim}{\to} H^*_{\D}(Y_y, F)$ $(F \in \Ob\, \D(Y))$, ce qui résulte du théorème 1.\footnote[30]{Donc on utilise \textbf{Der 5} (d).} (Cf. démonstration du théorème page 56, qui marche avec des hypothèses plus générales.)

On a c) $\Leftrightarrow$ c$'$) $\Rightarrow$ c$''$) (l'équivalence c) $\Leftrightarrow$ c$'$) par \textbf{Der 2}).

Prouvons c$''$) $\Rightarrow$ a).

\msp{page 70}

\noindent Considérons le diagramme à carrés cartésiens
\[
\begin{tikzcd}
X_{y_0} \arrow[r, "q"] \arrow[d, "p"{swap}] & X/y_0 \arrow[r] \arrow[d, "f_0"] & X \arrow[d, "f"] \\
\{y_0\} \arrow[r] & Y/y_0 \arrow[r] & Y
\end{tikzcd}
\qquad\qquad .
\]

\noindent L'hypothèse c$''$) prouve que pour tout $F$ dans $\D(X/y_0)$, on a $H^*_{\D}(X/y_0, F) \overset{\sim}{\to} H^*_{\D}(X_{y_0}, F)$. Par le théorème 1, ceci implique que $f$ est un $\D$-équivalence co-stricte. % [sic] un D-équivalence -> une D-équivalence, accord du genre

\noindent\textbf{Corollaire 1.} Énoncé dual du théorème, pour les foncteurs $X_{y_0} \to y_0\backslash X$, $X/x_0 \to Y/y_0$, et le changement de base des $j_!(F)$ \ldots

\noindent\textbf{Définition 4.} Sous les conditions équivalentes du théorème 1, on dit que $f$ est $\D$-\textit{propre.} La condition duale s'appelle $f$ $\D$-\textit{lisse.}

\noindent\textbf{Corollaire 2.} \textit{Soit} $f : X \to Y$ \textit{dans} $\Dia$. \textit{Si} $f$ \textit{est} $\D$\textit{-propre ou} $\D$\textit{-lisse, et à fibres} $\D$\textit{-contractiles, alors} $f$ \textit{est une} $\D$\textit{-équivalence stricte et costricte. Mieux,} $f$ \textit{est une} $\D$\textit{-équivalence} universelle \textit{i.e. est une} $\D$\textit{-équivalence et le reste après tout changement de base.}\footnote[31]{Se généralise : \textbf{Corollaire 2$'$.} \textit{Soit}
\[
\begin{tikzcd}[ampersand replacement=\&]
X \arrow[rr, "f"] \arrow[dr] \& \& Y \arrow[dl] \\
\& S \&
\end{tikzcd}
\]
\textit{avec} $X, Y$, $\D$\textit{-propres sur} $S$. \textit{Si} $f$ \textit{induit une} $\D$\textit{-équivalence sur les fibres c'est une} $\D$\textit{-équivalence. Itou dans le cas lisse.}\par
C'est la proposition 6 ci-dessous.}

\printedpage{42}
% head: A. Grothendieck : Dérivateurs, Ch. I, Édition des Universités de Montpellier II et Paris VII
% foot: Version 8 juillet 2001
\msp{page 71}

\noindent Il suffit de prouver que $f$ est une $\D$-équivalence. Supposons par exemple $f$ $\D$-propre. Il suffit de prouver que si $Y$ est $\D$-contractile, alors $X$ l'est (quitte à l'appliquer aux $Y/y_0$ après changement de base, pour montrer que les $X/y_0$ sont contractiles). Pour ceci, il suffit de prouver que $f$ est une $\D$-équivalence, et on sait que ceci est le cas si pour tout $F$ \textit{constant} sur $Y$, l'homomorphisme $F \overset{\text{adj}}{\longrightarrow} f_*f^*F$ est \texttt{[un]} isomorphisme. En fait, c'est vrai pour \textit{tout} $F$.

\noindent\textbf{Corollaire 3.} \textit{Si} $f : X \to Y$ \textit{est} $\D$\textit{-propre ou} $\D$\textit{-lisse, à fibres} $\D$\textit{-contractiles, alors} $f^* : \D(Y) \to \D(X)$ \textit{est pleinement fidèle, i.e. les homomorphismes d'adjonction}
\[
F \longrightarrow f_*f^*F \quad , \qquad f_!f^*F \longrightarrow F
\]
\textit{dans} $\D(Y)$ \textit{sont des isomorphismes.}

\noindent Il suffit de le voir pour $f$ propre, et cela résulte du calcul des fibres de $f_*f^*F$ \ldots

\noindent\textbf{Proposition 5.} \textit{La notion de foncteur} $\D$\textit{-propre} (resp. $\D$\textit{-lisse}) \textit{est stable par changement de base,} (trivial), \textit{et par composition} (utiliser \texttt{[le]} corollaire 2). \textit{Toute immersion fermée est} $\D$\textit{-propre, toute immersion ouverte} $\D$\textit{-lisse. Plus généralement, tout isomorphisme local est} $\D$\textit{-lisse, tout co-isomorphisme local est} $\D$\textit{-propre.}

\noindent\textbf{Proposition 6.} \textit{Soient} $X, Y$ \textit{sur} $S$, $f : X \to Y$ \textit{un} $S$\textit{-morphisme. Supposons} $X, Y$ \textit{propres sur} $S$, \textit{ou} $X$ \textit{et} $Y$ \textit{lisses sur} $S$, \textit{et que pour tout} $s \in \Ob\, S$, $f_s : X_s \to Y_s$ \textit{soit une} $\D$\textit{-équivalence, alors} $f$ \textit{aussi.}\footnote[32]{\textsc{Démonstration}. On prouvera même que sur la sous-catégorie des coefficients sur $Y$ qui proviennent de $\D(Y)$ \texttt{[plutôt de }$\D(S)$\texttt{]}, $f^* : \D(Y) \to \D(X)$ induit un foncteur pleinement fidèle. Si $f$ \texttt{[est]} propre, je l'interprète comme $q_*q^*F \overset{\sim}{\to} p_*p^*F$ pour tout $F$, et cela résulte du calcul de $p_*$, $q_*$ par passage aux fibres (propre), et de l'hypothèse sur les $f_s$. OK.}

\msp{page 72}

\noindent\textbf{Théorème 3.} \textit{Soit}
\[
\begin{tikzcd}
X \arrow[d, "f"{swap}] & X' \arrow[l, "q"{swap}] \arrow[d, "f'"] \\
Y & Y' \arrow[l, "p"{swap}]
\end{tikzcd}
\]
\textit{cartésien dans} $\Dia$. \textit{Alors le morphisme de changement de base}
\[
p^*f_* \longrightarrow f'_*q^*
\]
\textit{est un isomorphisme si} $f$ \textit{est} $\D$\textit{-propre, ou} $p$ $\D$\textit{-lisse.}

\noindent Le cas $\D$-propre déjà prouvé. Le cas $\D$-lisse : la démonstration p. 56 est OK.

\printedpage{43}
% head: A. Grothendieck : Dérivateurs, Ch. I, Édition des Universités de Montpellier II et Paris VII
% foot: Version 8 juillet 2001

\subsection*{Comparaison de $X$ et $X^{\circ}$.}

Soit $X$ \texttt{[une]} catégorie, on définit une catégorie $\Phi(X)$ et des foncteurs
\[
\begin{tikzcd}[column sep=small]
& & \Phi(X) \arrow[ddl, "p_X"{swap, name=P}] \arrow[ddr, "q_X"{name=Q}] & & \\
|[alias=L]| \text{\footnotesize\begin{tabular}{c}$\Cat$-cofibrant\\(donc propre)\end{tabular}} & & & & |[alias=R]| \text{\footnotesize\begin{tabular}{c}$\Cat$-fibrant\\(donc lisse)\\ \texttt{[plutôt }$\Cat$\texttt{-cofibrant}\\ \texttt{(donc propre)]}\end{tabular}} \\
& X & & X^{\circ} &
\arrow[from=P, to=L]
\arrow[from=Q, to=R]
\end{tikzcd}
\]
qui seront des $\hot$-équivalences, et des $\D$-équivalences si $\D$ est un dérivateur dont le domaine contient $X$ et $\Phi(X)$. On pose
\[
\Ob\, \Phi(X) = \Fl(X) \quad ,
\]
donc mêmes objets que $\mathcal{F}\ell(X) = \Hom(\boldsymbol{\Delta}_1, X)$, mais les morphismes sont différents :

\msp{page 73}
\[
\begin{tikzcd}
Y_1 \arrow[r, "u"] & Y_0 \arrow[d, "v_0"] \\
Y'_1 \arrow[u, "v_1"] \arrow[r, "u'"{swap}] & Y'_0
\end{tikzcd}
\]
un morphisme de $u$ dans $u'$ est un couple $v_0 : Y_0 \to Y'_0$, $v_1 : Y'_1 \to Y_1$ tel que $v_0uv_1 = u'$. Ainsi source et but $s(u)$ et $b(u)$ sont l'une contravariante, l'autre covariant en $u$. On prend pour $q$ le foncteur source, pour $p$ le foncteur but. % [sic] l'autre covariant -> l'autre covariante, accord avec « l'une contravariante »

Si on applique ceci à $X^{\circ}$, on trouve
\[
\begin{tikzcd}
& \Phi(X^{\circ}) \arrow[dl, "p_{X^{\circ}}"{swap}] \arrow[dr, "q_{X^{\circ}}"] & \\
X^{\circ} & & (X^{\circ})^{\circ} = X
\end{tikzcd}
\]
On a $\Ob\, \Phi(X^{\circ}) = \Fl(X^{\circ}) \simeq \Fl(X)$
\[
\begin{tikzcd}
Y_1 \arrow[r, "u"] & Y_0 \arrow[d, "v_1"] & \text{dans}\ \ X \\
Y'_1 \arrow[u, "v_0"] \arrow[r, "u'"{swap}] & Y'_0 & \text{dans}\ \ X
\end{tikzcd}
\]
et mêmes homomorphismes au sens de $\Phi(X)$ et au sens de $\Phi(X^{\circ})$. Donc
\[
\left\{
\begin{aligned}
&\Phi(X^{\circ}) = \Phi(X) \\
&p_{X^{\circ}} = q_X \\
&q_{X^{\circ}} = p_X \quad ,
\end{aligned}
\right.
\]

\end{document}
